\section{The symbols of positive and negative infinity}
\label{sec:positive-negative-infinity}

The symbols $+\infty$ and $-\infty$ record a failure of a real bound.
They are not real numbers and are not elements of any of the subsets of
$\R$ considered here. We use the extended ordering
\[
 -\infty<r<+\infty\qquad(r\in\R).
\]
For a nonempty set $E\subseteq\R$, the conventions in this section are
\[
 \begin{array}{ll}
 \sup E\in\R,&\text{if $E$ is bounded above},\\
 \sup E=+\infty,&\text{if $E$ is not bounded above},\\[2pt]
 \inf E\in\R,&\text{if $E$ is bounded below},\\
 \inf E=-\infty,&\text{if $E$ is not bounded below}.
 \end{array}
\]
The finite supremum exists by completeness, and the finite infimum
exists by its lower-bound version, Corollary~4.5. A set can have a finite
infimum while being unbounded above, or a finite supremum while being
unbounded below. A conclusion about one side does not require a bound
on the other side.

In explicit terms, not being bounded above means that for every real
$M$ there is an $e\in E$ with $e>M$. Not being bounded below means that
for every real $L$ there is an $e\in E$ with $e<L$. These statements
are the negations of the definitions of an upper bound and a lower bound.

For negation we use $-(+\infty)=-\infty$ and $-(-\infty)=+\infty$.
The additions needed in Exercise~5.7 are only
\[
 c+(-\infty)=(-\infty)+c=-\infty\quad(c\in\R),
 \qquad (-\infty)+(-\infty)=-\infty.
\]
These are extended-bound conventions. They do not make infinity a real
number, and they do not define the indeterminate expression
$(+\infty)+(-\infty)$.

\step{Two elementary algebra facts used for the interval descriptions}
First, if $0\le u<v$, then
\[
 v^2-u^2=(v-u)(v+u)>0.
\]
Both factors on the right are positive. Thus squaring preserves strict
order for nonnegative numbers. The corresponding weak inequality follows
by allowing $u=v$.

Second, cubing preserves strict order on all of $\R$. To check the
possible negative signs, suppose $u<v$. Then
\[
 v^3-u^3=(v-u)(v^2+vu+u^2),
\]
and
\[
 v^2+vu+u^2=\left(v+\frac u2\right)^2+\frac{3u^2}{4}.
\]
If $u\ne0$, the second term is positive and the first is nonnegative.
If $u=0$, then $v>0$ and the expression is $v^2>0$. Hence in both cases
$v^3-u^3>0$, because $v-u>0$. It follows, in particular, that
\[
 x^3\le8\ \Longleftrightarrow\ x\le2,
 \qquad
 x^3<8\ \Longleftrightarrow\ x<2.
\]
For example, $x\le2$ implies $x^3\le2^3=8$; conversely, $x^3\le8$
and $x>2$ would imply $x^3>8$, a contradiction. For the strict
equivalence, $x<2$ implies $x^3<8$, while $x\ge2$ would imply
$x^3\ge8$.

\exercise{5.1}{Four descriptions in interval notation}
\goal Write each of the following sets as an interval:
\[
 \begin{array}{ll}
 \text{(a)}\ \{x\in\R:x<0\},&
 \text{(b)}\ \{x\in\R:x^3\le8\},\\
 \text{(c)}\ \{x^2:x\in\R\},&
 \text{(d)}\ \{x\in\R:x^2<8\}.
 \end{array}
\]

\idea Translate the defining inequalities, and prove both inclusions
when the description involves squares. An infinite endpoint always has
a parenthesis, because it is not a real element that could be included.

\step{1. Part (a): All negative real numbers}
The notation $(-\infty,0)$ means precisely the real numbers strictly
less than $0$. If $x$ belongs to the given set, it satisfies $x<0$ and
therefore belongs to $(-\infty,0)$. Conversely, if $x\in(-\infty,0)$,
then $x\in\R$ and $x<0$, so it belongs to the given set. Thus
\[
 \{x\in\R:x<0\}=(-\infty,0).
\]
The endpoint $0$ is excluded because the defining inequality is strict.

\step{2. Part (b): A non-strict cubic inequality}
The cube-order calculation above gives both directions of
$x^3\le8\Longleftrightarrow x\le2$. Consequently
\[
 \{x\in\R:x^3\le8\}=\{x\in\R:x\le2\}=(-\infty,2].
\]
In particular, $2$ is included: $2^3=8$ satisfies the non-strict
inequality. This is why the finite endpoint has a square bracket.

\step{3. A completeness justification for the square roots needed below}
To identify the set of all squares, we need that every nonnegative real
number has a nonnegative real square root. Here is a proof using the
completeness axiom from Section~4.

The number $0$ has square root $0$. Fix $y>0$ and define
\[
 C=\{t\in\R:t\ge0\text{ and }t^2\le y\}.
\]
The set is nonempty because $0\in C$. It is bounded above by $y+1$:
if $t>y+1$, then $t>1$ and
$t^2>t>y$, so such a $t$ cannot belong to $C$. Completeness supplies
the real number $c=\sup C$, and $c\ge0$ since $0\in C$.

We claim $c^2=y$. If $c^2<y$, put $d=y-c^2>0$ and choose
\[
 h=\min\left\{1,\frac{d}{2(2c+1)}\right\}>0.
\]
We have $h\le1$ and $(2c+1)h\le d/2$. Expanding and using
$h^2\le h$ gives
\[
 (c+h)^2=c^2+2ch+h^2
 \le c^2+(2c+1)h
 \le c^2+\frac d2
 <c^2+d=y.
\]
Thus $c+h\in C$, but $c+h>c$, contradicting that $c$ is an upper
bound for $C$.

If $c^2>y$, then $c>0$. Put $d=c^2-y>0$ and choose
\[
 h=\min\left\{\frac c2,\frac{d}{4c}\right\}>0.
\]
Now $c-h\ge c/2>0$, and $2ch\le d/2$. Therefore
\[
 (c-h)^2=c^2-2ch+h^2
 \ge c^2-2ch
 \ge c^2-\frac d2
 >c^2-d=y.
\]
If $t\in C$ and $t>c-h$, both numbers are nonnegative, so the
square-order fact gives $t^2>(c-h)^2>y$, contradicting $t\in C$.
Thus $t\le c-h$ for every $t\in C$. The smaller number $c-h<c$
would be an upper bound for $C$, contradicting the least-upper-bound
property of $c$.

Both alternatives are impossible, so $c^2=y$. This proves existence
of a nonnegative square root. It is unique: if $a,b\ge0$ and
$a^2=b^2$, then $(a-b)(a+b)=0$. If $a+b=0$, nonnegativity forces
$a=b=0$; otherwise $a-b=0$, so again $a=b$. This justifies the
notation $\sqrt y$ without using any later result about continuity.

\step{4. Part (c): The set of all squares}
Let $P=\{x^2:x\in\R\}$. If $z\in P$, then $z=x^2$ for some
$x\in\R$. Every real square is nonnegative, so $z\ge0$ and
$z\in[0,+\infty)$. This proves $P\subseteq[0,+\infty)$.

Conversely, let $z\in[0,+\infty)$. Then $z\ge0$, and Step~3 gives
a real number $x=\sqrt z$ with $x^2=z$. Hence $z\in P$. We have
proved the reverse inclusion and therefore
\[
 \{x^2:x\in\R\}=[0,+\infty).
\]
The finite endpoint is included because $0=0^2$.

\step{5. Part (d): A strict square inequality}
Put $a=\sqrt8$. Step~3 gives $a\ge0$ and $a^2=8$; in fact $a>0$
because $8>0$. Let $D=\{x\in\R:x^2<8\}$.

First suppose $x\in D$. If $x\ge a$, nonnegative square order would
give $x^2\ge a^2=8$, a contradiction. Hence $x<a$. If $x\le-a$,
then $-x\ge a>0$, so $x^2=(-x)^2\ge a^2=8$, again a
contradiction. Thus $x>-a$. We have shown $-a<x<a$, so
$D\subseteq(-a,a)$.

Conversely, suppose $-a<x<a$. If $x\ge0$, then $0\le x<a$ and
$x^2<a^2=8$. If $x<0$, the inequality $x>-a$ becomes $-x<a$
after multiplying by $-1$, while $-x>0$. Therefore
$x^2=(-x)^2<a^2=8$ in this case too. Thus $x\in D$, proving the
reverse inclusion. We conclude
\[
 \{x\in\R:x^2<8\}=(-\sqrt8,\sqrt8).
\]
Neither endpoint belongs to the set, since
$(\sqrt8)^2=(-\sqrt8)^2=8$.

\answer{The intervals are
\textup{(a)} $(-\infty,0)$,
\textup{(b)} $(-\infty,2]$,
\textup{(c)} $[0,+\infty)$, and
\textup{(d)} $(-\sqrt8,\sqrt8)$.}

\exercise{5.2}{Extended bounds of the four intervals}
\goal Find the infimum and supremum, allowing the infinite symbols,
of each set in Exercise~5.1.

\idea A finite included endpoint is a minimum or maximum. For an open
endpoint, verify that elements approach it from the allowed side.
For an infinite endpoint, show explicitly that no real bound exists
on that side.

\step{1. Part (a): The interval $(-\infty,0)$}
The number $0$ is an upper bound, since every member is negative.
If $U<0$, let $x=U/2$. Then $x<0$, and
$x-U=-U/2>0$, so $x>U$. Thus no number less than $0$ is an upper
bound. This proves $\sup(-\infty,0)=0$.

For any proposed real lower bound $L$, put
$x=\min\{L,0\}-1$. Then $x<0$ and $x<L$, so $x$ belongs to the
interval and disproves that $L$ is a lower bound. The interval is not
bounded below, and its extended infimum is $-\infty$.

\step{2. Part (b): The interval $(-\infty,2]$}
Every element is at most $2$, and $2$ itself belongs to the interval.
Hence $2$ is its maximum. Every upper bound must be at least this
member $2$, so the supremum is $2$.

Given any real $L$, the number $x=\min\{L,2\}-1$ satisfies $x<2$
and $x<L$. It belongs to $(-\infty,2]$ and lies below the proposed
lower bound. Thus the interval is not bounded below and has infimum
$-\infty$.

\step{3. Part (c): The interval $[0,+\infty)$}
Every element is at least $0$, and $0$ belongs to the interval.
Thus $0$ is a minimum. Any lower bound is at most the member $0$,
so the greatest lower bound is $0$.

For any real proposed upper bound $M$, choose
$z=\max\{M,0\}+1$. Then $z\ge0$, so $z$ belongs to the interval,
and $z>M$. There is no real upper bound. Therefore its supremum is
$+\infty$.

\step{4. Part (d): The interval $(-\sqrt8,\sqrt8)$}
Write $a=\sqrt8>0$. Every $x\in(-a,a)$ satisfies $-a<x<a$.
Thus $-a$ is a lower bound and $a$ is an upper bound. We verify
the endpoint approximation required for them to be the greatest
lower bound and least upper bound.

Fix any $\varepsilon>0$ and set
\[
 \delta=\min\left\{\frac{\varepsilon}{2},a\right\}>0.
\]
Because $\delta\le a$, the numbers
\[
 x_-=-a+\delta,\qquad x_+=a-\delta
\]
both belong to $(-a,a)$: we have $-a<x_-\le0<a$ and
$-a<0\le x_+<a$. Moreover, $\delta\le\varepsilon/2<\varepsilon$,
so
\[
 x_-<-a+\varepsilon,
 \qquad x_+>a-\varepsilon.
\]
These inequalities show that no number greater than $-a$ can be
a lower bound, and no number less than $a$ can be an upper bound.
For example, if $L>-a$, choose $\varepsilon=L+a>0$ to obtain
$x_-<L$; if $U<a$, choose $\varepsilon=a-U>0$ to obtain $x_+>U$.
Therefore the infimum is $-a$ and the supremum is $a$.

\answer{The results, in the order of Exercise~5.1, are
\[
 \begin{array}{c|cc}
 \text{Part}&\inf&\sup\\ \hline
 \text{(a)}&-\infty&0\\
 \text{(b)}&-\infty&2\\
 \text{(c)}&0&+\infty\\
 \text{(d)}&-\sqrt8&\sqrt8
 \end{array}
\]
The finite endpoints in part \textup{(d)} are not attained.}

\exercise{5.3}{The unbounded sets from Exercise 4.1}
\goal Find the extended infimum and supremum of every unbounded set
in Exercise~4.1.

\idea The unbounded cases are precisely \textup{(h)}, \textup{(k)},
\textup{(l)}, \textup{(o)}, \textup{(t)}, and \textup{(u)}.
For each, determine its finite endpoint when one exists and exhibit
elements beyond every proposed bound on the unbounded side.

\step{1. Identify all the relevant cases}
We retain the letters of Exercise~4.1:
\begin{align*}
 H&=\bigcup_{n=1}^{\infty}[2n,2n+1],
 &K&=\left\{n+\frac{(-1)^n}{n}:n\in\N\right\},\\
 L&=\{r\in\Q:r<2\},
 &O&=\{x\in\R:x<0\},\\
 T&=\{x\in\R:x^3<8\},
 &P&=\{x^2:x\in\R\}.
\end{align*}
Here $P$ denotes the set in part \textup{(u)}. In particular, the
inequality in part \textup{(t)} is strict; it differs from the
non-strict cubic inequality in Exercise~5.1(b).

To check that no unbounded case was omitted, the remaining sets can
be enclosed in finite intervals as follows. The notation in the left
column refers to the letters of Exercise~4.1:
\begin{center}
\begin{tabular}{@{}ll@{}}
\toprule
Parts & A finite containing interval\\
\midrule
\textup{(a), (b), (e), (j), (s)} & $[0,1]$\\
\textup{(c)} & $[2,7]$\\
\textup{(f)} & $[0,0]$\\
\textup{(g)} & $[0,3]$\\
\textup{(i)} & $[-1,2]$\\
\textup{(m), (n)} & $[-2,2]$\\
\textup{(q)} & $[0,16]$\\
\textup{(r)} & $[0,2]$\\
\textup{(v), (w)} & $[-1,1]$\\
\bottomrule
\end{tabular}
\end{center}
For the reciprocal terms, $n\ge1$ gives $0<1/n\le1$, and
$0<1/3^n\le1$, which justify the bounds for \textup{(e), (j), (s)}.
Each intersection in \textup{(i)} and \textup{(r)} is a subset of
its $n=1$ interval, giving the displayed bounds. For \textup{(m)}
and \textup{(n)}, a real number outside $[-2,2]$ has square greater
than $4$ and therefore cannot have square less than $4$ or less than
$2$. The usual ranges of sine and cosine give the bounds for
\textup{(v)} and \textup{(w)}. The other displayed finite-interval
inclusions follow directly from their listed elements or intervals.
Finally, \textup{(d)} and \textup{(p)} are finite sets of real numbers.
For a finite list $a_1,\ldots,a_m$, the number
$C=1+\abs{a_1}+\cdots+\abs{a_m}$ satisfies
$-C<a_j<C$ for every listed element, so each such set is bounded.
This accounts for all the other parts.

\step{2. Part (h): An infinite union of intervals}
If $x\in H$, some $n\in\N$ satisfies $x\in[2n,2n+1]$.
Hence $x\ge2n\ge2$. Also $2\in[2,3]\subseteq H$, so $2$ is a
minimum and $\inf H=2$.

Given any real $M$, the Archimedean property supplies an integer
$n>\max\{1,M/2\}$. Then $2n>M$ and
$2n\in[2n,2n+1]\subseteq H$. Thus no real $M$ is an upper bound.
Therefore $\sup H=+\infty$.

\step{3. Part (k): A growing sequence with an alternating correction}
For every $n\in\N$, the value $(-1)^n$ is either $1$ or $-1$.
As $n>0$, it follows that
\[
 n+\frac{(-1)^n}{n}
 \ge n-\frac1n
 =\frac{n^2-1}{n}\ge0.
\]
The last inequality holds because $n\ge1$ gives $n^2-1\ge0$.
At $n=1$, the value is $1+(-1)/1=0$. Thus $0$ is the minimum
of $K$, and $\inf K=0$.

Also $1/n\le1$, so
\[
 n+\frac{(-1)^n}{n}\ge n-\frac1n\ge n-1.
\]
For a proposed real upper bound $M$, choose
$n>\max\{1,M+1\}$ by the Archimedean property. Then $n-1>M$,
and the corresponding member of $K$ is greater than $M$. Hence
$K$ is not bounded above and $\sup K=+\infty$.

\step{4. Part (l): The rational numbers less than $2$}
Every member of $L$ is less than $2$, so $2$ is an upper bound.
If $U<2$, density of $\Q$ in $\R$ supplies a rational number $r$
with $U<r<2$. This $r$ belongs to $L$ and exceeds $U$, so $U$
cannot be an upper bound. Thus $\sup L=2$, although $2\notin L$.

Given any real $D$, choose $n\in\N$ with
$n>\max\{1,-D\}$. Then $-n<D$, and the integer $-n$ is rational
and satisfies $-n<0<2$. Hence $-n\in L$ lies below $D$.
There is no real lower bound, so $\inf L=-\infty$.

\step{5. Part (o): The negative half-line}
This is $O=(-\infty,0)$, the set already considered in
Exercise~5.2(a). For completeness, $0$ bounds it above, and if
$U<0$ then $U/2\in O$ and $U/2>U$. Thus its supremum is $0$.
For every real $D$, the member $\min\{D,0\}-1$ lies below $D$.
Thus $O$ is not bounded below and $\inf O=-\infty$.

\step{6. Part (t): A strict cubic inequality}
By the cube-order fact proved at the start of the section,
\[
 T=\{x\in\R:x^3<8\}=(-\infty,2).
\]
All its elements are less than $2$, so $2$ is an upper bound.
If $U<2$, the number $x=(U+2)/2$ satisfies
\[
 x-U=\frac{2-U}{2}>0,
 \qquad 2-x=\frac{2-U}{2}>0.
\]
Hence $U<x<2$, so $x\in T$ and $U$ is not an upper bound.
Therefore $\sup T=2$; the supremum is not attained.

For any real $D$, set $x=\min\{D,2\}-1$. Then $x<2$ and $x<D$,
so $x\in T$ lies below $D$. This proves $\inf T=-\infty$.

\step{7. Part (u): The real squares}
Every real square is nonnegative and $0=0^2$ belongs to $P$.
Thus $\inf P=0$. Exercise~5.1(c) established, using completeness,
that $P=[0,+\infty)$. For any real $M$, the number
$y=\max\{M,0\}+1$ is nonnegative and exceeds $M$; since
$y=(\sqrt y)^2$, it belongs to $P$. No real upper bound exists,
and $\sup P=+\infty$.

\answer{Exactly six sets in Exercise~4.1 are unbounded. Their extended
bounds are
\[
 \begin{array}{c|cc}
 \text{Part}&\inf&\sup\\ \hline
 \text{(h)}&2&+\infty\\
 \text{(k)}&0&+\infty\\
 \text{(l)}&-\infty&2\\
 \text{(o)}&-\infty&0\\
 \text{(t)}&-\infty&2\\
 \text{(u)}&0&+\infty
 \end{array}
\]}

\exercise{5.4}{Reflection and extended bounds}
\goal For every nonempty $S\subseteq\R$, with
$-S=\{-s:s\in S\}$, prove
\[
 \inf S=-\sup(-S).
\]

\idea Negating an inequality changes a lower bound into an upper
bound. Treat the case of a finite infimum and the case of no real
lower bound separately.

\step{1. Establish how reflection changes bounds}
The set $-S$ is nonempty: if $s_0\in S$, then $-s_0\in-S$.
If a real number $d$ is a lower bound for $S$, then
\[
 d\le s\quad(s\in S)
 \quad\Longrightarrow\quad -s\le-d\quad(s\in S).
\]
Thus $-d$ is an upper bound for $-S$. Conversely, if $u$ is a real
upper bound for $-S$, then $-s\le u$ for all $s\in S$.
Multiplication by $-1$ reverses the inequality, giving $-u\le s$
for all $s\in S$. Hence $-u$ is a lower bound for $S$.
In particular, $S$ is bounded below if and only if $-S$ is bounded above.

\step{2. The case in which $S$ is bounded below}
Let $\alpha=\inf S\in\R$, whose existence follows from
Corollary~4.5. This step requires no upper bound on $S$.
Since $\alpha\le s$ for all $s\in S$, we have
$-s\le-\alpha$ for all $s\in S$. Therefore $-\alpha$ is an
upper bound for $-S$.

Now let $u$ be any real upper bound for $-S$. By Step~1, $-u$ is
a lower bound for $S$. The fact that $\alpha$ is the greatest lower
bound gives $-u\le\alpha$. Reversing signs gives
$-\alpha\le u$. Thus $-\alpha$ is less than or equal to every
upper bound for $-S$, proving
\[
 \sup(-S)=-\alpha.
\]
Negating this real equality yields
$-\sup(-S)=\alpha=\inf S$. This is the finite-infimum argument of
Exercise~4.9, with the bound correspondence written out explicitly.

\step{3. The case in which $S$ is not bounded below}
Here $\inf S=-\infty$ by definition. Let $M$ be any real number.
Since $S$ has no real lower bound, there exists $s\in S$ with
$s<-M$. Then $-s>M$, and $-s\in-S$. Thus $-S$ has elements
greater than every real $M$, so $\sup(-S)=+\infty$.
Using the stated negation convention gives
\[
 -\sup(-S)=-(+\infty)=-\infty=\inf S.
\]
These two cases exhaust a nonempty subset of $\R$, because its
infimum is either a real number or $-\infty$.

\answer{$\inf S=-\sup(-S)$ holds both for finite infima and for
the extended value $-\infty$.}

\exercise{5.5}{The infimum cannot exceed the supremum}
\goal Prove $\inf S\le\sup S$ for every nonempty $S\subseteq\R$,
whether or not $S$ is bounded on either side.

\idea Choose one real member of $S$. Its position between the
extended lower and upper bounds supplies a common comparison point.

\step{1. Compare the infimum to a chosen member}
Because $S$ is nonempty, choose $s_0\in S$. If $S$ is bounded below,
its finite infimum is a lower bound and therefore $\inf S\le s_0$.
If $S$ is not bounded below, then $\inf S=-\infty<s_0$ by the
extended order. Thus in either case
\[
 \inf S\le s_0.
\]

\step{2. Compare the same member to the supremum}
If $S$ is bounded above, its finite supremum is an upper bound and
therefore $s_0\le\sup S$. If $S$ is not bounded above, then
$\sup S=+\infty>s_0$. Hence in either case
\[
 s_0\le\sup S.
\]

\step{3. Use transitivity of the extended order}
Combining the preceding inequalities gives
\[
 \inf S\le s_0\le\sup S.
\]
The desired inequality follows. This argument includes the cases in
which one or both bounds are infinite, without treating infinity as
an element of $S$.

\answer{For every nonempty $S\subseteq\R$, $\inf S\le\sup S$.}

\exercise{5.6}{A chain of bounds under set inclusion}
\goal If $S,T$ are nonempty subsets of $\R$ and $S\subseteq T$,
prove
\[
 \inf T\le\inf S\le\sup S\le\sup T.
\]

\idea A lower bound for the larger set is also a lower bound for the
smaller set; the same holds for upper bounds. Check the infinite cases
on the relevant side before applying the finite greatest- or
least-bound property.

\step{1. Prove the leftmost inequality}
If $T$ is not bounded below, then $\inf T=-\infty$.
The number $\inf S$ is either real or $-\infty$, so the extended
order gives $\inf T\le\inf S$, including equality when both are
$-\infty$.

If $T$ is bounded below, put $\alpha=\inf T\in\R$.
Every $t\in T$ satisfies $\alpha\le t$. Because every $s\in S$
also belongs to $T$, we have $\alpha\le s$ for all $s\in S$.
Thus $S$ is bounded below as well, and $\alpha$ is one of its real
lower bounds. Its greatest lower bound must be at least this lower
bound, so
\[
 \inf T=\alpha\le\inf S.
\]
This argument assumes only a lower bound for $T$, not an upper bound.

\step{2. Prove the middle inequality}
Exercise~5.5 applies to the nonempty set $S$ and gives
\[
 \inf S\le\sup S.
\]
In particular, this remains valid if $S$ is unbounded on one or both
sides.

\step{3. Prove the rightmost inequality}
If $T$ is not bounded above, then $\sup T=+\infty$.
Since $\sup S$ is either real or $+\infty$, we have
$\sup S\le\sup T$ in the extended order, including equality when
both are $+\infty$.

If $T$ is bounded above, put $\beta=\sup T\in\R$.
Every element of $T$ is at most $\beta$, so every element of the
subset $S$ is at most $\beta$. Thus $S$ is bounded above, and
$\beta$ is an upper bound for it. The least-upper-bound property gives
\[
 \sup S\le\beta=\sup T.
\]
Here only an upper bound for $T$ was needed.

\step{4. Assemble the comparisons}
The three inequalities have been proved in all finite and infinite
cases. Placing them in order yields the required chain.

\answer{For nonempty $S\subseteq T\subseteq\R$,
$\inf T\le\inf S\le\sup S\le\sup T$.}

\exercise{5.7}{The infimum of a sum set, including unbounded cases}
\goal For arbitrary nonempty $A,B\subseteq\R$, define
\[
 A+B=\{a+b:a\in A,\ b\in B\}.
\]
Prove the extended equality
\[
 \inf(A+B)=\inf A+\inf B.
\]

\idea When both sets are bounded below, add their lower bounds and
use the defining approximation property of an infimum with an
$\varepsilon/2$ allowance for each summand. When a set is unbounded
below, fix one member of the other set and show that the sum set is
unbounded below too.

\step{1. Check nonemptiness and the possible extended values}
Choose $a_0\in A$ and $b_0\in B$, which is possible because both
sets are nonempty. The real number $a_0+b_0$ belongs to $A+B$,
so the sum set is nonempty. Put
\[
 \alpha=\inf A,\qquad \beta=\inf B.
\]
Each of $\alpha,\beta$ is either a real number or $-\infty$.
Neither can be $+\infty$ under the conventions for nonempty subsets
of $\R$. Thus the expression $\alpha+\beta$ never involves the
undefined combination $+\infty+(-\infty)$.

\step{2. Suppose both $A$ and $B$ are bounded below}
In this case $\alpha,\beta\in\R$. For every $a\in A$ and $b\in B$,
the lower-bound properties give $a\ge\alpha$ and $b\ge\beta$.
Adding yields
\[
 a+b\ge\alpha+\beta.
\]
Hence $\alpha+\beta$ is a real lower bound for the nonempty set
$A+B$. Corollary~4.5 supplies its real infimum
$\gamma=\inf(A+B)$. Because $\gamma$ is the greatest lower bound,
\[
 \gamma\ge\alpha+\beta.
\]

To prove the reverse inequality, fix any $\varepsilon>0$.
The real number $\alpha+\varepsilon/2$ is larger than the greatest
lower bound $\alpha$, so it cannot be a lower bound for $A$.
Consequently some $a\in A$ satisfies
\[
 \alpha\le a<\alpha+\frac{\varepsilon}{2}.
\]
The first inequality holds because $\alpha$ is a lower bound.
The same reasoning for $B$ supplies $b\in B$ with
\[
 \beta\le b<\beta+\frac{\varepsilon}{2}.
\]
Adding the two strict upper estimates gives
\[
 a+b<\alpha+\beta+\varepsilon.
\]
Since $a+b\in A+B$ and $\gamma$ is a lower bound for that set,
\[
 \gamma\le a+b<\alpha+\beta+\varepsilon.
\]
This holds for every $\varepsilon>0$. If
$\gamma>\alpha+\beta$, choose the positive real number
$\varepsilon=\gamma-(\alpha+\beta)$. The displayed inequality
would then give
\[
 \gamma<\alpha+\beta+
 \bigl(\gamma-(\alpha+\beta)\bigr)=\gamma,
\]
a contradiction. Therefore $\gamma\le\alpha+\beta$.
Combining both comparisons proves
\[
 \inf(A+B)=\alpha+\beta
\]
when both infima are finite. No upper bound on $A$ or $B$ was used.
Thus this case also covers sets that are unbounded above but bounded
below.

\step{3. Suppose $A$ is not bounded below}
Here $\alpha=-\infty$. Keep the fixed real member $b_0\in B$ from
Step~1. Let $D$ be any proposed real lower bound for $A+B$.
Because $A$ is not bounded below and $D-b_0$ is real, there is an
$a\in A$ satisfying
\[
 a<D-b_0.
\]
Adding the real number $b_0$ gives $a+b_0<D$. The number
$a+b_0$ belongs to $A+B$, so $D$ is not a lower bound.
As this works for every real $D$, the sum set is unbounded below,
and $\inf(A+B)=-\infty$.

If $B$ is bounded below, then $\beta$ is finite and
$\alpha+\beta=(-\infty)+\beta=-\infty$.
If $B$ is also unbounded below, then $\beta=-\infty$ and
$\alpha+\beta=(-\infty)+(-\infty)=-\infty$.
Thus the equality holds in both possibilities for $B$.

\step{4. Suppose $A$ is bounded below and $B$ is not}
The case with both sets unbounded below was already covered in Step~3.
It remains to handle a finite $\alpha$ and $\beta=-\infty$.
Fix the real member $a_0\in A$. Given any real $D$, the absence
of a lower bound for $B$ supplies $b\in B$ with $b<D-a_0$.
Then $a_0+b<D$ and $a_0+b\in A+B$. Hence $A+B$ is unbounded
below and
\[
 \inf(A+B)=-\infty
 =\alpha+(-\infty)
 =\inf A+\inf B.
\]

\step{5. Account for all cases}
Either both sets are bounded below, or $A$ is unbounded below, or
$A$ is bounded below and $B$ is unbounded below. Steps~2--4 prove
the equality in these exhaustive cases. Every sum used to produce a
member of $A+B$ was a sum of two real numbers; the infinite symbols
were used only to record the resulting extended bounds.

\answer{For all nonempty $A,B\subseteq\R$,
$\inf(A+B)=\inf A+\inf B$, with the stated extended-bound conventions.}
